\section{A positive statistic and its binary-tree transfers}
\label{sec:construction}

We combine the size and collision estimates with the Fourier supply
of Proposition~\ref{prop:supply} to construct a nonnegative statistic.
Poisson summation will turn it into an amplitude.  The transfers
propagate a quantitative lower bound once their diagonals are controlled;
Section~\ref{sec:conclusion} will bound the same amplitude from above
by averaging over arrangements of its prime variables.

The pivot extraction, extension to positive integers, and doubling of
the remaining template follow the architecture of
Section~\ref{sec:characters}.  The exact residue-mask transforms carry
no prescribed multiplicative-character phase, so the anchor cancellation
is replaced by a common outside list of spectator primes, where
Lemma~\ref{lem:tree-comparison} compares the paired histories.

All parameters introduced in this section are new; in particular, the
letters $L,k,m$ no longer denote parameters from the preceding character
arguments.

\subsection{Scales and prime priors}

Let $L\to\infty$.  The transfer depth $k$ will be a sufficiently large
fixed positive integer, chosen before $L$ tends to infinity.  Put
\[
 h=e^{.01L},\qquad z=k^4,\qquad
 m=2\lfloor zL/2\rfloor,\qquad r_j=2^{j-1}\quad(1\le j\le k).
\]
The notation $o_k(m)$ and $O_k(1)$ permits dependence on all parameters
that are fixed before $L$.  In bounds of the form $Cm$ or $C2^lm$,
however, $C$ will be independent of $k$ and of the large gap constants
introduced below, once $k$ and then $L$ are sufficiently large.  The
threshold on $L$ may depend on these fixed choices.

Choose fixed positive constants $B_s,B_D,B_z$, with their order of
choice specified in Section~\ref{sec:conclusion}, and define
\begin{align}
 \Delta_0&=(B_s+8\log z)m,\notag\\
 \Delta_j&=r_j(B_D+B_z\log z+.4\log r_j)m
       \qquad(1\le j\le k),\notag\\
 E_l&=\Delta_0+\sum_{j=1}^l\Delta_j+2^l\sqrt m,\notag\\
 V_l&=e^{E_l}\qquad(0\le l\le k).
 \label{eq:src43}
\end{align}
Thus $\log V_l=O_k(m)$.  In particular, all frequencies used below
are smaller than every sampled prime, since even the smallest prime
range has lower endpoint $\exp(\exp(.0005L))$.

Fix a nonnegative smooth function $\varphi$ supported on $[-1,1]$
such that
\[
 \sum_{n\in\Z}\varphi(t-n)=1\qquad(t\in\R).
\]
Then $0\le\varphi\le1$ and $\int_{\R}\varphi(t)\,dt=1$.
For a center $c$, the \emph{log-cell prior} is the probability measure
on primes proportional to
\[
                 \frac{\varphi(\log p-c)}p.
\]
Except for the giant-prime prior defined shortly, all prime priors
will omit a set $\mathcal E_L$ of at most two prime values.  This set
is fixed for each $L$; Section~\ref{sec:arithmetic} specifies its
choice from the two possible Page exceptional conductors.  Every
abundance assertion in this section holds uniformly after any such
deletion.  A nongiant cell is therefore normalized by
\[
 Z_c=\sum_{\substack{p\ \text{prime}\\p\notin\mathcal E_L}}
                 \frac{\varphi(\log p-c)}p.
\]
For all the cell centers used here, the prime number theorem and
partial summation give $Z_c\sim c^{-1}$.  The same statement holds
without the deletion.  They also give
\[
 \sum_p\frac{\log p}{p}\varphi(\log p-c)=1+o(1).
\]
The errors are uniform as the centers in our specified ranges tend
to infinity.  Deleting two prime values does not change these
asymptotics.

By \eqref{eq:src28}, there is a block $[G_0,G_0+h]$ in log-prime
coordinates, contained in the range of that estimate, such that the
primes satisfying its two conditions have total $(\log p)/p$ weight
at least $c_0h$, for a fixed $c_0>0$.  Here and below a prime satisfying
those conditions is called favorable.  To see the block assertion,
partition
$[e^{.05L},e^{.9L}]$ into intervals of length $h$, apart from the two
end pieces.  On each full block the log coordinate varies by a
relative $o(1)$, uniformly in the block.  If every block had favorable
$(\log p)/p$ weight at most $ch$, summing these bounds after division
by the left log coordinate would give favorable harmonic weight at
most $(.85c+o(1))L$.  A sufficiently small fixed $c$ contradicts
\eqref{eq:src28}.  Mertens' estimates make the two end pieces
negligible in this calculation.

Set
\[
                   \log X=2G_0+2^{k+1}h,
       \qquad A'=A\cap[X^{9/10},X].
\]
The whole chosen block lies below the prime cutoff in
\eqref{eq:src2} for the uniform measure on $A'$.  Indeed,
$\log\sqrt X=G_0+2^kh$, whereas the block ends at $G_0+h$, and the
loss from the logarithmic denominator in that cutoff is only
$O(\log\log X)=O(L)$.

Use the transforms $F_p,g_p$ from \eqref{eq:src27}, with the
normalization
\[
 g_p(b)=\sqrt p\,\E_{x\bmod p}F_p(x)e_p(-bx),\qquad
 F_p(x)=p^{-1/2}\sum_{b\bmod p}g_p(b)e_p(bx).
\]
For a favorable prime put
\[
 \widetilde g_p(b)=
 \begin{cases}
   g_p(b)/|g_p(b)|,&g_p(b)\ne0,\\
   0,&g_p(b)=0,
 \end{cases}
\]
and put $\widetilde g_p=0$ at every other prime.  Let
$\widetilde F_p$ be its inverse transform with the same
normalization.  Conjugate symmetry of $g_p$ implies that
$\widetilde F_p$ is real.  Also
\[
 \E_x\widetilde F_p(x)=0,
 \qquad \E_x|\widetilde F_p(x)|^2
      =\E_b|\widetilde g_p(b)|^2\le1.
\]
For a favorable prime, Parseval and the definition of $F_p$ give
\[
 \E_{x\in S_p}\widetilde F_p(x)
   =\sqrt{\frac{1-\sigma_p}{\sigma_p}}\,
        \E_b|g_p(b)|\ge\frac{\delta_0}{\sqrt2}.
\]

There is an integer center $G$, whose log-cell support is contained
in $[G_0,G_0+h]$, for which
\begin{equation}
 \E_{\substack{p\ \text{in the cell at}\ G\\a\in A'}}
                     \widetilde F_p(a)\ge c_1>0.
 \label{eq:src44}
\end{equation}
The giant prior in this display includes all primes in its cell.
Here is why the collision estimate applies to these possibly
unbounded inverse transforms.  If $\mu_p$ is the projection of the
uniform probability on $A'$, then Cauchy--Schwarz gives
\[
 \left|\sum_{x\bmod p}
    (\mu_p(x)-U_{S_p}(x))\widetilde F_p(x)\right|^2
       \le p\,\|\mu_p-U_{S_p}\|_2^2.
\]
Consequently \eqref{eq:src2} bounds the sum of these squared errors
with weight $(\log p)/p$ by $O(L)$.  The total such weight on the
block is $O(h)$, so its weighted sum of absolute errors is
$O(\sqrt{Lh})=o(h)$.  The favorable uniform means, on the other
hand, have a positive sum of order $h$.  Partitioning by the integer
translates of $\varphi$, and deleting the bounded-width end strips,
therefore supplies a cell with a positive average empirical mean.
Replacing $(\log p)/p$ by $1/p$ in this cell does not change this
conclusion: the log coordinate has relative variation $o(1)$, and
the same error bound applies.  This proves \eqref{eq:src44}.

Fix this $G$.  Write
\begin{equation}
 Z_G=\sum_{p\ \text{prime}}\frac{\varphi(\log p-G)}p,
 \qquad
 \mu_G(p)=\frac{\varphi(\log p-G)}{pZ_G},
 \qquad c_g=\log Z_G^{-1}.
 \label{eq:construction-giant-prior}
\end{equation}
We have $Z_G\sim G^{-1}$, and hence
$c_g=\log G+O(1)\le .91L$ for sufficiently large $L$.
The coefficient in this last bound is independent of the gap
constants and of $k$.

\subsection{The initial lists and positive statistic}

A \emph{half-list} consists of the following independent positions:
\begin{enumerate}[label=(\roman*)]
 \item one giant with prior $\mu_G$, using the test
       $\widetilde F_p$;
 \item $m/2$ bulk positions, with their common prior proportional to
       $1/p$ on
       $.004L\le\log\log p\le .006L$,
       omitting $\mathcal E_L$;
 \item $m/2$ spectator positions, with their common prior proportional
       to $1/p$ on a subset of
       $.0005L\le\log\log p\le .001L$
       of harmonic mass $\gg L$, omitting $\mathcal E_L$, on which
       \eqref{eq:src27} holds uniformly with a bound tending to zero;
 \item three top positions and two compensation positions of each
       type $j=1,\ldots,k$, with cell priors chosen below.
\end{enumerate}
Every nongiant position uses the exact test $F_p$.  The spectator
subset exists by Corollary~\ref{cor:mixed-flatness}; its exceptional
threshold can be chosen slowly enough that all the stated
conclusions hold uniformly on the retained subset.  The bulk and
spectator harmonic normalizing masses are both between fixed
positive multiples of $L$.

Let $S_b,S_d$ be the sums of the log primes at the bulk and spectator
positions of one half-list.  Include in its weight the two factors
\[
                   \varphi(S_b-t_b)\varphi(S_d-t_d).
\]
These are weights, not conditioning of any prior.  Choose integer
centers $t_b,t_d$ so that, for each of the two groups, the expectation
of its cutoff together with the restriction
$1/3\le\sigma_p\le2/3$ at all its positions is at least $e^{-Cm}$.
Such choices follow directly from the partition of unity.  Almost
all the probability of a single position is balanced, by
\eqref{eq:src27} or \eqref{eq:src2}; for sufficiently large $L$ the
probability that all positions in that group are balanced is at
least $2^{-m/2}$.  Each possible log sum lies in an interval admitting
at most $\exp(O(L)+O(\log m))$ relevant integer centers.  At least
one center thus has the required weighted mass.  In particular,
\[
                    t_b,t_d=o_k(h).
\]

Call the bulk and top positions \emph{protected}.  For the full list
of two half-lists, define the target log sum $J$ for the protected
positions and the target log sum $w_j$ for compensation type $j$ by
\begin{equation}
 J=\frac{\log X-2G-2t_d+\Delta_0+\sum_{j=1}^k\Delta_j}{2^k},
 \qquad
 w_j=J+\sum_{i>j}w_i-\frac{\Delta_j}{r_j}.
 \label{eq:src45}
\end{equation}
The second definition is made successively for $j=k,k-1,\ldots,1$.
It gives
\[
 J\asymp h,\qquad
 w_j=(2^{k-j}+o_k(1))J,
 \qquad
 J+\sum_{j=1}^kw_j=2^kJ-\sum_{j=1}^k\Delta_j.
\]
Choose the three top cell centers in a fixed small relative neighborhood of
$(J-2t_b)/6$, with their sum equal to $(J-2t_b)/2+O(1)$.
For the two positions of type $j$, choose centers in a corresponding
neighborhood of $w_j/4$, with sum $w_j/2+O(1)$.  Neighborhoods of
relative radius $.03$ suffice to keep different compensation types,
top positions, bulk positions, spectator positions, and giants in
pairwise disjoint prime ranges.

All these top and compensation centers can be chosen so that the
balanced primes have a fixed positive fraction of the cell prior.
Indeed, these primes also lie below the cutoff in
\eqref{eq:src2}.  A cell in which a fixed positive fraction is
unbalanced contributes a fixed positive amount to the nonnegative
imbalance term in that estimate.  Since the translates of
$\varphi$ sum to one, only $O(L)$ integer centers can fail the desired
balanced-fraction condition.  Each neighborhood available here has
length comparable to $h$, which is much larger than $L$.
For a prescribed rounded pair sum, exclude a bad center and its
reflection about that sum; this excludes only $O(L)$ choices.
For a triple, choose the first two centers in slightly smaller
neighborhoods.  Among the resulting pairs only an $O(L/h)$ fraction
have a bad third center determined by the rounded sum.  This gives
the required centers while respecting all ranges.  On the bulk
cutoff, the protected log sum of each half-list is now
$J/2+O(1)$.

Let $T(n)$ be the expectation over one half-list of the product of
all its tests at $n$, multiplied by the two log-sum cutoffs.  For
$a\in A'$, every exact local test is
\[
                     F_p(a)=\sqrt{\frac{1-\sigma_p}{\sigma_p}}>0.
\]
This value is independent of the choice of $a\in A'$: all these
primes are sufficiently small for the defining residue restriction
on $A$ to apply.  The expectation of the product of exact tests and
the cutoffs is therefore a positive scalar $c_L$, independent of
$a$, with $c_L\ge e^{-Cm}$.  This lower bound follows by restricting
every nongiant position to balanced primes; each local value is at
least $2^{-1/2}$, and the preceding bin and cell choices give the
required mass.  Thus
\[
                    T(a)=c_L\E_{p\sim\mu_G}\widetilde F_p(a).
\]
Take the nonnegative Schwartz function $\psi$ used in the quadratic
argument, with compactly supported smooth Fourier transform and a
positive lower bound on $[0,1]$.  The tests are real, so $T(n)^2$ is
nonnegative.  Equation~\eqref{eq:src44}, Jensen's inequality on $A'$,
and the lower bound in \eqref{eq:src1} give
\begin{equation}
              I:=\sum_{n\in\Z}\psi(n/X)T(n)^2
                   \ge \sqrt X\,e^{-Cm}.
 \label{eq:src46}
\end{equation}
The logarithmic loss in $|A'|$ is absorbed here because
$\log\log X=O(L)$ and $m\asymp k^4L$.  Notice that this argument
uses the mean of the giant test and its square; it does not require
the giant test to be pointwise nonnegative.

\subsection{Removing repeated primes and applying Poisson summation}

Expand $T(n)^2$ using two independent half-lists.  Let $M$ be the
product of all positions, counted with multiplicity.  On the cutoff
support,
\[
 \log M=2G+2t_d+J+\sum_{j=1}^kw_j+O_k(1)
        =\log X+\Delta_0+O_k(1).
\]
We may discard the tuples with repeated prime values, with a
negligible error relative to \eqref{eq:src46}.  We give the estimate
because no cancellation estimate for repeated local factors is
being assumed.

For a tuple containing a repetition, the product of its distinct
primes, denoted $Q$, satisfies
\[
 Q\ll_k X e^{\Delta_0}/\exp(\exp(.0005L))=o(X).
\]
The product of the local tests is periodic modulo $Q$.  Compact
Fourier support of $\widehat\psi$ implies that its smoothed sum is
exactly $X\widehat\psi(0)$ times its residue mean, for sufficiently
large $L$.  If any prime occurs once, that residue mean is zero by
the mean-zero property of its test and CRT.  If a prime occurs with
multiplicity $e\ge2$, its contribution to the absolute residue
mean is at most $p^{e/2-1}$.  Indeed every local test has probability
$L^2$ norm at most one, hence supremum at most $\sqrt p$; apply
Cauchy--Schwarz to two factors and use the supremum bound for the
others.  This argument also covers factors assigned different
roles at the same prime.

The probability of any prescribed ordered tuple is at most
\[
            M^{-1}L^{-2m}\exp(Cm+O_k(1)).
\]
There are $2m$ broad-band positions.  Each contributes a
normalization of size at most $C/L$.  The remaining $4k+8$ cell
positions have combined normalization at most $\exp(O(kL)+O_k(1))$,
which is $\exp(Cm+O_k(1))$ with $C$ independent of large $k$ because
$m\asymp k^4L$.  Since
$\prod_{p\mid Q}p^{e_p/2-1}=M^{1/2}/Q$, the absolute contribution
of tuples in which every prime occurs at least twice is at most
\[
 \sqrt X\exp(-\Delta_0/2+Cm+O_k(1))L^{-2m}
               \sum_{\rm tuples}\frac1{\prod_{p\ \text{distinct}}p}.
\]
Put $b=2m+4k+8$.  The harmonic mass of the union of all prime ranges
is $O(L)$.  Assigning the $b$ labeled positions to $i$ distinct
values, and then dropping restrictions for an upper bound, gives
\[
 \sum_{\rm tuples}\frac1{\prod_{p\ \text{distinct}}p}
    \le \sum_{i=1}^{b}\frac{i^b(CL)^i}{i!}
    \le b^b e^{CL}.
\]
Thus the repeated-prime contribution is bounded by
\[
 \sqrt X\exp\left(-\frac{\Delta_0}{2}
                +(2\log z+C)m+o_k(m)\right).
\]
Choosing $B_s$ sufficiently large makes this negligible in
\eqref{eq:src46}.  The constant required here is independent of the
later gap constants and of sufficiently large $k$.

Retain only tuples with all positions distinct, and let $\eta_0$ be
their Poisson expression divided by $\sqrt X$.  We have shown that
\begin{equation}
                        |\eta_0|\ge e^{-B_*m}
 \label{eq:src47}
\end{equation}
for a fixed $B_*$ independent of $B_D,B_z$ and $k$.  The same
$B_*$ may be used for all sufficiently large admissible $B_s$.

\subsection{Templates, coefficients, and support conventions}
\label{subsec:coefficient-support}

Keep the $m$ spectator positions of the two half-lists, with their
original half-list roles, as one outside list $\mathbf d$, and write
$d$ for their product.  This list will never be duplicated.  All
other positions form the ordered regular template $\mathcal I_0$;
its two giants are labeled $+$ and $-$ in a fixed order.

At step $j$, split $\mathcal I_{j-1}$ into the pivot positions and
the remaining template $H$.  The pivot product is $P=pu$, where
$p$ is the $+$ giant and $\mathbf u$ is the ordered list of all
current compensation positions of type $j$, with product $u$.
Unlike the whole-pivot extension in Section~\ref{sec:characters},
here only the giant $p$ will be extended to positive integers.
The entries of $\mathbf u$ retain their actual-prime priors.
Thus $H$ contains the other giant, all protected positions, and all
positions of types $i>j$.  Form
$\mathcal I_j=(H_+,H_-)$ from two copies of $H$, labeling their
giants $+$ and $-$, respectively.  Whenever a template or a part of
one is sampled, its positions have independent copies of their
specified priors.

Before step $j$, there are $r_j$ copies of the full protected list
of two half-lists, and $r_j$ copies of each compensation type not
yet removed.  In particular $\mathbf u$ has $4r_j$ positions.
On the inherited protected bins,
\begin{align}
 \log u&=r_jw_j+O_k(1),\notag\\
 \log H&=T_H+O_k(1),\qquad
 T_H=G+r_j\left(J+\sum_{i>j}w_i\right)
       =G+r_jw_j+\Delta_j.
 \label{eq:src48}
\end{align}
Here a template and its product are denoted by the same letter only
when no confusion can result.

If $R$ is the product of distinct actual regular prime values, put
\[
 G_R(b)=\prod_{\ell\mid R}
             g_\ell^{\rm reg}\bigl(b/(R/\ell)\bmod\ell\bigr),
\]
where $g_\ell^{\rm reg}=\widetilde g_\ell$ on a giant position and
$g_\ell^{\rm reg}=g_\ell$ on every small regular position.  Field
fractions always have unit denominators.  This expression is used
only on the pairwise coprime support, and is extended by zero when
that support fails.  The amplitudes will have the form
\begin{equation}
 \eta_l=\E_{\mathbf d,\mathcal I_l}
       \sum_{0<|s|\le V_l}G_R(s/d)A_l(s),
                  \qquad R=\prod\mathcal I_l.
 \label{eq:src49}
\end{equation}
The factor $G_R(s/d)$ depends on the actual regular prime values and
their product, so it is unchanged when the bulk values are reassigned
among their positions.  The coefficient $A_l$ will contain the
spectator factors and all arrangement-dependent history weights and
support conditions.  This separates the common factor from the
coefficient to be averaged over arrangements.
The outer positions in this expectation are actual primes.
The coefficient $A_l$, unlike $G_R$, will also be defined when its
two giant entries are positive integers.

We specify all its support conventions.  A history is a full binary
tree of depth $l$, with leaves at level zero.  Each node at level
$i$ has a nonzero signed integer frequency of magnitude at most
$V_i$.  At every state in a history require pairwise coprimality of
all its current slots and the outside list $\mathbf d$.  Its two
current giant values must also be units against \emph{every}
frequency at that node or below it.  These last restrictions may
always be added at an actual top, without changing an expression:
actual giant primes exceed every frequency in any of the finitely
many possible history levels.  Once added, they are kept when a
giant is extended to an integer.  These requirements are imposed
termwise on complete histories; they include cross-branch
frequency tests for either current giant.  A newly inserted pivot
is a current giant in its child states, so it is subject to their
subtree tests.  Failed support conditions give zero.

At level zero set
\begin{equation}
 \begin{split}
 A_0(s)={}&\left(\frac{X}{dR}\right)^{1/2}
             \widehat\psi\left(-\frac{sX}{dR}\right)\,
       \prod_{\rm halves}\varphi(S_b-t_b)\varphi(S_d-t_d)\\
 &\hspace{13mm}\cdot
       \prod_{q\mid d}
          g_q\bigl(s/((d/q)R)\bmod q\bigr),
 \end{split}
 \label{eq:src50}
\end{equation}
on this support, and zero otherwise.  The half-list spectator
weights use the original outside roles of $\mathbf d$.
Formula~\eqref{eq:src49} with $l=0$ is the distinct-tuple Poisson
formula.  Indeed the local inverse expansions contribute
$(dR)^{-1/2}$, and Poisson summation followed by division by
$\sqrt X$ gives the factor $\sqrt{X/(dR)}$ in
\eqref{eq:src50}.  The zero frequency vanishes since the local
transforms vanish at zero.  On the cutoffs,
$dR/X=\exp(\Delta_0+O_k(1))$; the additional $\sqrt m$ in
$E_0$ therefore accommodates the fixed compact support of
$\widehat\psi$.

For the recursive definition of $A_j$, a current state is
$\mathcal I_j=(H_+,H_-)$ with root frequency $s$.
Sum over $v,w$ with $0<|v|,|w|\le V_{j-1}$, and independently sample
the type-$j$ vector $\mathbf u$ by its priors.  Restore the pivot by
\begin{equation}
                       p=\frac{vH_- - wH_+}{su}.
 \label{eq:src51}
\end{equation}
Retain it only if it is a positive integer.  The two child states
are copies of $\mathcal I_{j-1}$ with entries $(p,\mathbf u,H_+)$
and $(p,\mathbf u,H_-)$ in their prescribed roles.  Their giant
pairs, in order, are the inserted $p$ and the giant of $H_+$ or
$H_-$, respectively.  Keep the same outside list $\mathbf d$.
Multiply the two child coefficients, conjugating the right one,
by
\begin{equation}
                         u\varphi(\log p-G).
 \label{eq:src52}
\end{equation}
Sum and average these terms, imposing all the history restrictions
above.  In a formula, with state arguments displayed only here,
\begin{align*}
 A_j(s;\mathbf d,H_+,H_-)
   =\sum_{\substack{0<|v|\le V_{j-1}\\0<|w|\le V_{j-1}}}
       \E_{\mathbf u}\bigl[{}
       &u\varphi(\log p-G)
          A_{j-1}(v;\mathbf d,p,\mathbf u,H_+)\\
       &\cdot\overline{A_{j-1}(w;\mathbf d,p,\mathbf u,H_-)}
                              \bigr].
\end{align*}
The bracket is interpreted termwise after expanding the two child
histories, and is zero unless \eqref{eq:src51} and all their support
conditions hold.  Internal samples in the two children are
independent; only $\mathbf u$ at the present node is shared.
This defines $A_l$ for actual or integer giant entries, and includes
no transform attached to an integer giant.

Figure~\ref{fig:transfer-recursion} records the dependencies in this recursion.

\begin{figure}[tbp]
  \centering
  \begin{tikzpicture}[
    every node/.style={font=\small,align=center},
    transfer box/.style={
      draw=black!65,rounded corners=2pt,
      inner xsep=6pt,inner ysep=5pt
    },
    transfer arrow/.style={-{Latex[length=2mm]},line width=.6pt}
  ]
    \node[transfer box,fill=black!3,text width=14.4cm] (current) {
      Current state $\mathcal I_j=(H_+,H_-)$, root frequency $s$\\[2pt]
      Sum over $0<|v|,|w|\le V_{j-1}$;
      sample one type-$j$ vector $\mathbf u$.\\[3pt]
      $\displaystyle p=\frac{vH_- - wH_+}{su},
        \qquad u=\prod_a u_a$\\[3pt]
      Node factor, used once: $u\varphi(\log p-G)$
    };

    \node[transfer box,text width=6.55cm,anchor=north]
      (left) at ($(current.south)+(-3.75cm,-.65cm)$) {
      Left child $(p,\mathbf u,H_+)$\\[3pt]
      $A_{j-1}(v;\mathbf d,p,\mathbf u,H_+)$
    };
    \node[transfer box,text width=6.55cm,anchor=north]
      (right) at ($(current.south)+(3.75cm,-.65cm)$) {
      Right child $(p,\mathbf u,H_-)$\\[3pt]
      $\overline{A_{j-1}(w;\mathbf d,p,\mathbf u,H_-)}$
    };

    \coordinate (split) at ($(current.south)+(0,-.25cm)$);
    \draw[line width=.6pt] (current.south) -- (split);
    \draw[transfer arrow] (split) -| (left.north);
    \draw[transfer arrow] (split) -| (right.north);

    \node[anchor=north,text width=14.8cm,inner sep=0pt]
      at ($(left.south)!.5!(right.south)+(0,-.28cm)$) {
      The same $p$, $\mathbf u$, and outside spectator list $\mathbf d$
      are used in both children.\\[2pt]
      Retain only $p\in\mathbb Z_{>0}$;
      all support zero conventions of Subsection~\ref{subsec:coefficient-support} apply.
    };
  \end{tikzpicture}
  \caption{The coefficient recursion for $A_j(s)$:
    multiply the node factor by the two displayed child factors,
    with no prime-specific transform attached to the integer pivot $p$.
    Subsequent internal samples in the two descendants are independent
    in their underlying priors, before the complete-history support
    indicators are applied.}
  \label{fig:transfer-recursion}
\end{figure}

\subsection{The exact transfer and its diagonal}

We verify the relation between the successive expressions
\eqref{eq:src49}.  Fix $\mathbf d,p,\mathbf u$ at step $j$, put
$P=pu$, and group old terms according to
$t=v/H\pmod P$.  On the original prime support,
\[
 G_{PH}(v/d)=G_P(t/d)G_H(v/(Pd)).
\]
Let
\[
 B_t=\E_H
       \sum_{\substack{0<|v|\le V_{j-1}\\v/H\equiv t\ ({\rm mod}\ P)}}
                G_H(v/(Pd))A_{j-1}(v),
\]
with all inherited support conditions; invalid terms are zero.
The extracted row has counting squared norm at most $P$:
\[
                   \sum_{t\bmod P}|G_P(t/d)|^2\le P.
\]
This is CRT and the exact probability squared norm one at each
compensation prime, together with squared norm at most one at the
giant.  Cauchy--Schwarz in this row, followed by Jensen over the
outer probability measures, yields
\[
                  |\eta_{j-1}|^2
                    \le\E_{\mathbf d,p,\mathbf u}
                                  P\sum_{t\bmod P}|B_t|^2.
\]
Crucially, $B_t$ contains no prime-specific transform of the pivot
giant $p$.  It is defined for a positive integer $p$ by exactly the
same formulas, with the same coprimality and frequency filters.
Equation~\eqref{eq:construction-giant-prior} gives
$p\mu_G(p)=e^{c_g}\varphi(\log p-G)$ at primes.  Extending the
nonnegative sum to all positive integers gives
\begin{equation}
 |\eta_{j-1}|^2
   \le e^{c_g}\E_{\mathbf d,\mathbf u}
          \sum_{p\in\N}u\varphi(\log p-G)
                         \sum_{t\bmod pu}|B_t|^2.
 \label{eq:src53}
\end{equation}
Here the sum over $p\in\N$ means positive integers.

Expand the square using independent $H_+,H_-$ and child frequencies
$v,w$.  Separate the terms $vH_-=wH_+$ as the diagonal.
For every other term, the congruence of the two group indices
defines the nonzero signed integer
$s=(vH_--wH_+)/(pu)$, so the substitution is exactly
\eqref{eq:src51}.  It is one-to-one, with $p$ still any positive
integer in its cell.  Equations~\eqref{eq:src43} and
\eqref{eq:src48} show
\[
 \log|s|\le E_{j-1}+\Delta_j+O_k(1)<E_j.
\]
The margin is $2^{j-1}\sqrt m$, which tends to infinity.
All protected bins used in this size estimate are present in each
nonzero child history.

The actual lists $H_+,H_-$ cannot share a prime on such a term.
If $\ell$ divided both products, it would divide $s$, since it is
coprime to $pu$; but $\ell>|s|$.  At a prime of $H_+$ and a prime
of $H_-$, respectively, \eqref{eq:src51} gives
\[
                 v/P=s/H_-,\qquad -w/P=s/H_+.
\]
The inverse tests are real, so
$g_\ell^{\rm reg}(-b)=\overline{g_\ell^{\rm reg}(b)}$.
The two regular transforms therefore combine exactly into
$G_{H_+H_-}(s/d)$.  The remaining weights, histories, and inserted
pivot are precisely the recursive definition of $A_j$.
Additional giant-frequency tests at the new top are free because
its two giants are actual primes.  Consequently the off-diagonal
part of the extended expression in \eqref{eq:src53}, before
$e^{c_g}$, is exactly $\eta_j$.

Write $\mathcal D_j$ for its exact diagonal.  It is nonnegative, as
we now verify, and we have
\begin{equation}
                e^{-c_g}|\eta_{j-1}|^2
                         \le\mathcal D_j+\eta_j.
 \label{eq:construction-transfer-split}
\end{equation}
In particular the off-diagonal expression $\eta_j$ is real: the
extended square and its diagonal are real.

On the diagonal, all prime factors of $H_+$ and $H_-$ exceed the
frequencies, so $vH_-=wH_+$ forces $H_+=H_-$ as products and
$v=w$.  Grouping by this common product and frequency \emph{before}
expanding the square gives a common factor
$|G_H(v/(pud))|^2$ times the squared absolute value of the
prior-weighted sum of coefficients over its ordered arrangements.
This proves nonnegativity.  The common factor is independent of the
ordering: each prime uses the same local transform in every
arrangement, and the single giant in $H$ is fixed by its disjoint
range.  We may therefore drop its bounded giant factor in this
nonnegative expression.  The remaining multiplier is
\begin{equation}
 \mathcal G=\prod_{\ell\mid H_{\rm sm}}
       \left|g_\ell\bigl(v/(pud(H/\ell))\bmod\ell\bigr)\right|^2,
 \label{eq:src54}
\end{equation}
where $H_{\rm sm}$ consists of the small slots of $H$.

Now expand by bijections matching the slots in the two $H$ lists.
This counts every ordered counterpart once, since the prime values
within a valid $H$ are distinct.  Only matches within the same
prime band can occur.  In particular the giant matches itself and
bulk positions match bulk positions.

Fix a first ordered $H$.  For any prescribed counterpart ordering,
its point probability is $C_H/H$ times its smooth cell factors and
its role-membership indicators.  The normalizing constants satisfy
\begin{equation}
                 C_H\le L^{-r_jm}\exp(Cr_jm+O_k(1)).
 \label{eq:construction-counterpart-mass}
\end{equation}
Indeed there are $r_jm$ bulk positions, and $O(kr_j)$ cell positions.
The latter normalizations cost $\exp(O(kr_jL)+O_k(1))$, absorbed
uniformly by $\exp(Cr_jm)$ because $m\asymp k^4L$.
On simultaneous coefficient support, \eqref{eq:src48} allows us to
extract $C_He^{-T_H}$, leaving $e^{T_H}/H=O_k(1)$.
Normalize the external integer measure as
\begin{equation}
       \int f(p)\,d\mu_{\mathrm{ext}}(p)
          :=e^{-G}\sum_{p\ge1}\varphi(\log p-G)f(p).
 \label{eq:construction-external-measure}
\end{equation}
It has bounded total mass.  Extracting $e^G$ from its definition
and $e^{r_jw_j}$ from $u$ leaves the net scalar
$C_He^{-\Delta_j}$.  The remaining Archimedean factor is
\[
 K_{\rm ar}
   =\frac{e^{T_H}}H\frac{u}{e^{r_jw_j}}
        \prod_{\text{counterpart cell positions}}
                           \varphi(\log p_i-c_i).
\]
This factor is $O_k(1)$ on simultaneous coefficient support.
Counterpart role-membership indicators are retained as zero
conventions in addition to these displayed smooth factors.

\subsection{The two comparison estimates}

We state precisely the comparison estimates needed for the diagonal
and the final symmetrization.  Their proof occupies
Section~\ref{sec:arithmetic}.  At level $l$, put $r=2^l$.
Each coefficient $A_l$ has $r$ bottom leaves, each carrying the
$m$ bulk positions of the initial full list.  Its $m$ outside
spectator positions are shared by all leaves.

There are two outer environments, always using the priors, integer
extension, coefficients, and support conventions just defined.
\begin{enumerate}[label=(\alph*)]
 \item In the \emph{final environment}, $l=k$, the current list is
       $\mathcal I_k$ with its actual prime priors.  Set
       $\mathcal G=1$.  A pair of assignments may permute all $rm$
       bulk values among their labeled positions; the nonbulk
       positions and the outside spectator list are common.
       Set $K_{\rm ar}=1$.
 \item In the \emph{diagonal environment} for step $j$, $l=j-1$.
       The current $+$ giant $p$ has the external measure
       \eqref{eq:construction-external-measure}.  All other current
       slots, including the type-$j$ vector $\mathbf u$, have their
       original actual prime priors.  Set $\mathcal G$ equal to
       \eqref{eq:src54} with $v=s$, with the zero convention on
       invalid top support.  In a pair of assignments the first
       uses the ordered $H$ and the second uses a prescribed
       same-band matching of its slots, with $p,\mathbf u,
       \mathbf d$ unchanged.  Use the factor $K_{\rm ar}$ above,
       including the counterpart role indicators.  For the norm
       of a \emph{single} assignment this factor is omitted.
\end{enumerate}
The expectation in the second environment includes the
bounded-mass integer measure; it need not be a probability measure.
All root-frequency sums below are over $0<|s|\le V_l$.

For a pair of assignments, form the bipartite multigraph whose two
vertex sets are their respective $r$ bottom leaves.  Each actual
bulk variable is an edge joining the leaf that contains it in the
first assignment to the leaf that contains it in the second.
Every vertex has degree $m$.  The graph and its number of connected
components depend only on the slot assignments, not on the sampled
prime values or on internal histories.

\begin{proposition}[Comparison estimates]
\label{prop:comparisons}
In either of the two outer environments, the following estimates
hold uniformly over the indicated assignments.  For one
assignment,
\begin{equation}
       \E\sum_{0<|s|\le V_l}\mathcal G|A_l(s)|^2
           \le\exp\bigl(r(\Delta_0+Cm)+o_k(m)\bigr).
 \label{eq:src55}
\end{equation}
The constant $C$ is independent of sufficiently large $k$ and of
the fixed gap budgets.  For a pair of assignments, if $l\ge2$ and
its overlap graph has at most $3r/4$ connected components, then
\begin{equation}
 \left|\E\sum_{0<|s|\le V_l}\mathcal G K_{\rm ar}
            A_l^{(1)}(s)\overline{A_l^{(2)}(s)}\right|
                        \le\exp(-\omega_k(m)).
 \label{eq:src56}
\end{equation}
Here $\omega_k(m)/m\to\infty$ as $L\to\infty$ for every fixed
choice of the other parameters.  Equivalently, the right side of
\eqref{eq:src56} is smaller than $\exp(-C_km)$ for every fixed
$C_k$ once $L$ is sufficiently large.  The estimate is uniform in
the allowed slot matchings and permutations.
\end{proposition}

The assumptions about all current-state coprimalities and all
giant-frequency unit tests are part of this proposition.  They are
used in its arithmetic reduction, not additional conclusions of
the estimate.  The outside spectator list is not resampled within
either coefficient or between a paired comparison.  In contrast,
the two coefficients have independent internal prime samples;
within each coefficient, a sample at one node is shared by its
two children as specified in \eqref{eq:src52}.  The assigned actual
top values have the stated coupling between the coefficients.
Coincidences among independently sampled internal prime values
are allowed when the support permits them, and their equality
patterns will be summed explicitly in the proof.
